\section{Conclusion}\label{sec:conclusion}

This paper presented \scl{}, a training-free layer that sits between an unchanged detector and an
unchanged tracker and adapts the tracker's input to the detector's scores, using only the scores of
recent frames and the tracks of the previous frame. Nested Otsu splits of the recent logits decide
whether the scores are clean or noisy. On clean scores the layer keeps the tracker's settings; on
noisy scores it chooses the candidates, gives them new scores, and sets the thresholds; and a
continuation rescue gives single-stage trackers the step they lack. A host contract of four
published settings lets one fixed configuration serve many trackers. The layer was developed step
by step from \aob{} and was then frozen before any external test.

One frozen configuration was tested on seven published trackers, four detection sources, and three
datasets. With halved detector scores, it restored two trackers from 0 to about 66 HOTA. It raised
the HOTA of both single-stage trackers, including the external PD-SORT, on every MOT17 sequence.
With RT-DETR-L on KITTI, it raised the MOTA of ByteTrack and BoT-SORT by 5.98 and 7.86 points and
cut their identity switches by 43\% and 47\%, and on held-out aerial data it raised MOTA by 24.70
points and cut identity switches by up to 56\%. It reached higher HOTA than simple adaptive
thresholds on the same scores and than the published ACT method under halved scores and on
held-out aerial data, kept well-matched two-stage trackers at their own accuracy, showed no
statistically significant HOTA loss in any tested cell, and took 2.95 ms per frame on a CPU.

Overall, the results show that the score scale of the detector should be treated as part of the
tracking system rather than as a fixed setting of the tracker. Future work will test the layer with
more detectors, trackers, and datasets.
